\chapter{Те же узлы на Verilog}
\label{ch:bridge}

В предыдущих главах вы строили схемы из блоков: рисовали вентили в Конструкторе,
собирали сумматор, запоминали тревогу в триггере и вели сквозной проект
«Охранная панель». Теперь то же самое говорится словами на языке Verilog.
Вопрос главы такой: как каждый знакомый узел записывается кодом и как
убедиться, что получилась именно та схема, которую вы собирали руками?

Ответ мы будем искать одним и тем же способом. Берём узел, о котором уже знаем
всё из теоретических глав, описываем его несколькими строками, нажимаем
«Синтезировать» и сравниваем вкладку RTL с собственной схемой. Диаграммы
подтверждают таблицу истинности или поведение во времени. В конце главы все
узлы сойдутся в «Охранной панели» на Verilog.

\begin{table}[H]
\centering
\caption{Узел из теоретических глав и его запись на Verilog}
\small
\begin{tabularx}{\textwidth}{l X X}
\toprule
\textbf{Узел} & \textbf{Схема из блоков} & \textbf{Строка на Verilog} \\
\midrule
Логика & И, ИЛИ, НЕ в Конструкторе & \code{assign alarm = door \& (night | armed);} \\
Сумматор & цепочка полных сумматоров & \code{assign \{co, s\} = a + b + ci;} \\
Сравнение & компаратор «A больше B» & \code{assign loud = noise > limit;} \\
Дешифратор & четыре элемента И & \code{assign y = en ? (4'd1 << a) : 4'd0;} \\
Триггер & D-триггер с тактом & \code{always @(posedge clk) q <= d;} \\
Счётчик & регистр и сумматор «+1» & \code{q <= q + 1;} \\
Автомат & регистр состояния и логика переходов & \code{case (state) ... endcase} \\
\bottomrule
\end{tabularx}
\end{table}

Исходники лежат в папках \code{code/00\_theory/} (узлы этой главы) и
\code{code/05a\_bridge/} (итоговый проект); примеры остальных глав~--- в
папках с именами глав (\code{code/b9a\_fsm/}, \code{code/05\_verilog/} и т.\,д.). Откройте любой файл в
Logra (проект «С нуля», вкладка «Код», группа «Создать»), нажмите
«Синтезировать (Ctrl+Enter)» и повторяйте за книгой. Результат смотрите во
вкладках группы «Проверить»: RTL, Диаграммы и Лаборатория.

\section{Сумматор: формулы превращаются в вентили}

Полный сумматор из главы~«Арифметические узлы» описывается двумя формулами.
В Verilog каждая формула --- одна строка \code{assign}.

\vfile{\codepath/00_theory/full_adder.v}{code/00\_theory/full\_adder.v}

После синтеза вкладка RTL показывает ровно то, что мы рисовали в первой
части: два элемента XOR, два AND и OR (рис.~\ref{fig:b-adder-rtl}). Каждый
вентиль подписан номером строки в исходном тексте, а на проводах стоят имена
сигналов.

\shot{br-adder-rtl}{RTL-схема полного сумматора: вентили из учебника, подписи берутся из кода}{fig:b-adder-rtl}

Проверим таблицу истинности. Кнопка «Перебрать входы» на вкладке «Диаграммы»
подаёт по очереди все восемь комбинаций $a$, $b$ и $c_{in}$, а на
диаграмме (рис.~\ref{fig:y-wave}) видно, как отвечают $sum$ и $c_{out}$.
Сравните с таблицей в главе~«Арифметические узлы»: $sum$ равна единице при
нечётном числе единиц на входах, а $c_{out}$ --- когда единиц две или три.

Теперь самое интересное. В главе~\ref{ch:inside} мы говорили, что любой логический
элемент в ПЛИС заменяется таблицей --- LUT. Откроем вкладку Gowin и щёлкнем
по LUT для $c_{out}$.

\shot{br-adder-lut}{Тот же сумматор на примитивах Gowin: две LUT3 и таблица $c_{out}$ в инспекторе}{fig:b-adder-lut}

Синтезатор заменил все пять вентилей \emph{двумя} таблицами (по одной на
выход). Конфигурация LUT для $c_{out}$ --- число \code{0xE8}, то есть
\code{11101000$_2$}: единицы стоят в позициях 3, 5, 6 и 7. Это адреса
\code{011}, \code{101}, \code{110} и \code{111}, то есть как раз комбинации,
в которых хотя бы две из трёх входных переменных равны единице.

\begin{ideabox}
Вентили из теоретических глав и таблицы LUT --- два способа записать одну и ту же
булеву функцию. На бумаге мы выбираем вентили, в ПЛИС синтезатор вычисляет
таблицу сам. Результат для нас один и тот же: таблица истинности.
\end{ideabox}

\begin{diybox}[Сумматор и порог шума]
\textbf{Задание.} Опишите 4-разрядный сумматор с переносом и сравнением и проверьте, что он считает так же, как цепочка из полных сумматоров в Конструкторе.

\textbf{В Logra.} Проект «С нуля», вкладка «Код». Введите \code{assign \{co, s\} = a + b + ci;} с \code{input [3:0] a, b}, \code{input ci}, \code{output [3:0] s}, \code{output co}, нажмите «Синтезировать». Для порога шума добавьте \code{output loud} и строку \code{assign loud = noise > 4'd9;} (вход \code{noise} на 4 бита).

\textbf{Ожидаемый результат.} На вкладке RTL --- сумматор и компаратор; $9+8+1$ даёт \code{s = 2} и \code{co = 1}; \code{loud} равен 0 при \code{noise = 9} и 1 при \code{noise = 10}.

\textbf{Как проверить.} Вкладка «Диаграммы»: задайте $a=9$, $b=8$, $ci=1$ и сверьте $s$ и $co$; затем перебор значений \code{noise}. На вкладке Gowin запишите, сколько ячеек занял сумматор.

\stars{1} повторите: полный сумматор \code{full\_adder.v} из этой главы, сценарий «Перебрать входы», сверка с таблицей. \stars{2} измените: сделайте вычитатель \code{assign \{bo, d\} = \{1'b0, a\} - \{1'b0, b\};} и найдите пару, при которой \code{bo = 1} (заём). \stars{3} проект: этап «Охранная панель» --- тревога по порогу (\code{loud}); критерий приёмки --- на диаграмме \code{loud} поднимается ровно при \code{noise} от 10 до 15.
\end{diybox}

\section{Дешифратор: сдвиг вместо перебора}

Дешифратор $2\to 4$ с входом разрешения можно описать и таблицей, и
формулами для каждого выхода. В Verilog проще всего воспользоваться сдвигом:
единицу сдвигаем влево на столько позиций, сколько говорит адрес.

\vfile{\codepath/00_theory/decoder2to4.v}{code/00\_theory/decoder2to4.v}

\shot{br-dec-rtl}{Дешифратор: сдвигатель и мультиплексор, выбирающий «всё нули» при $en = 0$}{fig:b-dec-rtl}

Среда нарисовала другую схему, чем наша в учебнике: вместо четырёх элементов
И она показала один сдвигатель \code{<<} и мультиплексор. Обе схемы делают
одно и то же. RTL показывает схему \emph{по смыслу} кода, а на кристалле её
всё равно заменят четыре LUT. Диаграмма на рис.~\ref{fig:b-dec-wave}
подтверждает: на любом адресе горит ровно одна линия, а при $en = 0$ --- ни
одной.

\shot{br-dec-wave}{Перебор всех 8 комбинаций входов дешифратора}{fig:b-dec-wave}

\begin{diybox}[Выбор зоны дешифратором]
\textbf{Задание.} Охранная панель следит за четырьмя зонами. Адрес зоны \code{z} (2 бита) должен зажигать ровно одну из четырёх ламп.

\textbf{В Logra.} Вкладка «Код»: откройте \code{decoder2to4.v} из этой главы (или введите свой модуль) и подключите выходы к четырём лампам зон. Нажмите «Синтезировать».

\textbf{Ожидаемый результат.} При $en=1$ горит одна лампа, номер которой равен адресу; при $en=0$ не горит ни одной.

\textbf{Как проверить.} Диаграммы: сценарий «Перебрать входы» даёт восемь комбинаций, на каждой не больше одной единицы на выходе. Вкладка RTL: найдите сдвигатель \code{<<}; сравните с вашей схемой из четырёх элементов И.

\stars{1} повторите: проверьте все 8 строк перебора. \stars{2} измените: сделайте дешифратор $3\to 8$ через сдвиг \code{8'd1 << a} и убедитесь, что схема стала шире, а запись --- не длиннее. \stars{3} проект: этап «Охранная панель» --- зона тревоги; критерий приёмки --- при срабатывании датчика зоны 2 горит только лампа 2, а при сбросе все лампы гаснут.
\end{diybox}

\section{Защёлка и триггер: два вида памяти рядом}

Чем защёлка отличается от триггера, мы подробно разбирали в главе
«Триггеры». Теперь это можно увидеть в схеме. Опишем оба элемента в одном
модуле: защёлку --- блоком \code{always @*} с условием по уровню, триггер ---
блоком \code{always @(posedge clk)}.

\vfile{\codepath/00_theory/dff_vs_latch.v}{code/00\_theory/dff\_vs\_latch.v}

\shot{br-latch-rtl}{Слева триггер DFF, который реагирует на фронт тактового сигнала $C$; справа защёлка LATCH с входом разрешения $EN$}{fig:b-latch-rtl}

Разница в условных обозначениях --- та же, что на рисунках глав~\ref{ch:ff}--\ref{ch:regs}:
у триггера на тактовом входе стоит треугольник (фронт), у защёлки вход
разрешения $EN$ действует по уровню.

\begin{warnbox}
Защёлка возникает в Verilog \emph{по ошибке} чаще, чем по замыслу: если в
комбинационном блоке не присвоить значение при каком-то условии, синтезатор
поставит защёлку, чтобы сохранить старое значение. Так что строку
\code{LATCH} на RTL-схеме, которую вы не планировали, стоит считать
сигналом тревоги.
\end{warnbox}

\begin{diybox}[Запомнить тревогу]
\textbf{Задание.} Тревога должна оставаться включённой, даже если дверь уже закрыли, пока её не снимет ключ. В главе~\ref{ch:ff} для этого был триггер; опишите его на Verilog.

\textbf{В Logra.} Вкладка «Код», новый модуль с входами \code{clk}, \code{trig}, \code{key}, выходом \code{mem}. Тело модуля: \code{always @(posedge clk) if (key) mem <= 1'b0; else if (trig) mem <= 1'b1;}

\textbf{Ожидаемый результат.} После импульса \code{trig} выход \code{mem} остаётся единицей, пока не придёт \code{key}; на RTL --- один триггер DFF с логикой сброса и установки.

\textbf{Как проверить.} Диаграммы: подайте \code{trig} на один такт, затем снимите его, \code{mem} должна держаться; потом подайте \code{key}. В RTL не должно быть ни одного элемента LATCH.

\stars{1} повторите: найдите DFF и LATCH в \code{dff\_vs\_latch.v}. \stars{2} найдите ошибку: уберите ветку \code{else} в блоке \code{always @*}, который выставляет \code{y} только при \code{en = 1}, и объясните, откуда на схеме взялась защёлка. \stars{3} проект: этап «Охранная панель» --- запоминание тревоги; критерий приёмки --- на диаграмме \code{mem} не отпускает тревогу после закрытия двери и сбрасывается только по \code{key}.
\end{diybox}

\section{Счётчик: регистр плюс сумматор}

Синхронный счётчик из главы «Регистры и счётчики» состоит из регистра и
схемы прибавления единицы. Verilog описывает это одной строкой
\code{q <= q + 1}, а сброс и разрешение счёта --- условными операторами.

\vfile{\codepath/00_theory/counter4.v}{code/00\_theory/counter4.v}

\shot{br-cnt-rtl}{RTL-схема счётчика: сумматор $+1$ и 4-разрядный регистр с синхронным сбросом и входом разрешения}{fig:b-cnt-rtl}

На схеме легко найти все части, о которых мы говорили: сумматор
(круг с плюсом), регистр (прямоугольник DFF), вход разрешения \code{EN},
синхронный сброс \code{SRST}. Выход регистра возвращается на вход
сумматора --- это и есть обратная связь, превращающая комбинационный
сумматор в счётчик.

Диаграмма (рис.~\ref{fig:b-cnt-wave}) показывает счёт. Сначала действует
сброс (\code{rst = 1}), потом счёт идёт по возрастающим фронтам такта: $0, 1,
2, \ldots, 15$, а затем $q$ возвращается к нулю --- модуль счёта равен 16.

\shot{br-cnt-wave}{Работа счётчика: сброс, счёт и переполнение}{fig:b-cnt-wave}

\begin{diybox}[Счётчик времени тревоги]
\textbf{Задание.} Охранная панель должна знать, сколько тактов длится тревога, и не переполняться: дойдя до 15, счётчик останавливается.

\textbf{В Logra.} Вкладка «Код». Возьмите \code{counter4.v} и замените условие счёта на \code{if (en \&\& q != 4'd15) q <= q + 4'd1;}. Синтезируйте и подключите \code{q} к четырём светодиодам.

\textbf{Ожидаемый результат.} При \code{en = 1} значение растёт $0,1,\dots,15$ и остаётся равным 15 (насыщение); в исходном \code{counter4} после 15 было бы 0.

\textbf{Как проверить.} Диаграммы: поставьте число циклов 20 и держите \code{en = 1}. На RTL найдите сумматор «+1» и регистр с входом разрешения, как на рис.~\ref{fig:b-cnt-rtl}.

\stars{1} повторите: запустите исходный \code{counter4} и прочитайте модуль счёта 16 по диаграмме. \stars{2} измените: сделайте счётчик по модулю 10 (после 9 снова 0). \stars{3} проект: этап «Охранная панель» --- время тревоги; критерий приёмки --- пока \code{state} равно ALARM, счётчик растёт до 15 и остаётся на 15, при сбросе возвращается к 0.
\end{diybox}

\section{Автоматы и очередь: что где разобрано}
\label{sec:b-fsm-ref}

Конечный автомат~--- это регистр состояния и логика переходов
(глава~\ref{ch:sync}). Автомат Мура для светофора имеет три состояния; в каждом
горит одна лампа, а переход определяется числом прошедших тактов. Граф вы уже
строили из блоков, здесь та же структура записывается кодом.

\begin{figure}[H]
\centering
\begin{tikzpicture}[st/.style={circle, draw=FpgaBlue, fill=FpgaLight, minimum size=18mm, font=\sffamily\small, align=center, thick}]
  \node[st] (g) at (0,0) {GREEN\\{\footnotesize 4 такта}};
  \node[st] (y) at (5,0) {YELLOW\\{\footnotesize 2 такта}};
  \node[st] (r) at (10,0) {RED\\{\footnotesize 4 такта}};
  \draw[arr] (g) -- node[above, lbl] {$t=3$} (y);
  \draw[arr] (y) -- node[above, lbl] {$t=1$} (r);
  \draw[arr] (r) to[bend left=35] node[below, lbl] {$t=3$} (g);
  \draw[arr] (-2.4,0) -- node[above, lbl] {rst} (g);
\end{tikzpicture}
\caption{Граф состояний светофора. $t$ --- счётчик тактов, проведённых в состоянии}
\label{fig:b-fsm-graph}
\end{figure}

Код этого автомата с тремя блоками (регистр состояния, переходы, выходы),
автомат по продаже (Мили), детектор последовательности 1-0-1, три способа
кодирования состояний и две очереди FIFO с признаками \code{full} и
\code{empty} разобраны с тестбенчами в главе~\ref{ch:fsmdesign}, в разделе
«Автомат в три блока на Verilog». Исходники лежат в \code{code/b9a\_fsm/} и
\code{code/00\_theory/} (там же прежний вариант светофора \code{traffic\_fsm.v}
и диаграмма его работы, рис.~\ref{fig:b-fsm-wave}). Автомат режимов панели из
главы~\ref{ch:fsmdesign} входит в итоговый проект этой главы.

\section{Код Хэмминга}

Кодер, канал с ошибкой и декодер из главы~\ref{ch:codes} целиком состоят из
исключающих ИЛИ. Для проверки в модуле предусмотрен вход \code{err}: маска,
которая переворачивает выбранные биты в «канале».

\vfile{\codepath/00_theory/hamming74.v}{code/00\_theory/hamming74.v}


Проверяющий модуль \code{tb\_hamming74.v} (тестбенч, глава~\ref{ch:sim}; лежит
рядом) перебирает все 16
значений данных и все семь одиночных ошибок и убеждается, что данные каждый
раз восстанавливаются и что синдром указывает на повреждённую позицию.

\begin{diybox}[Защита кода зоны от ошибки]
\textbf{Задание.} Номер зоны (4 бита) уходит по длинному проводу, и в нём может перевернуться один бит. Проверьте, что код Хэмминга исправляет любую одиночную ошибку.

\textbf{В Logra.} Вкладка «Код»: \code{hamming74.v}. Вход \code{d} --- ваш номер зоны, \code{err} --- маска повреждения. Тестбенч \code{tb\_hamming74.v} лежит рядом.

\textbf{Ожидаемый результат.} При \code{err = 0000100} выходы \code{syn} равны 3, \code{q} равно исходному \code{d}. При \code{err = 0000000} синдром 0.

\textbf{Как проверить.} Диаграммы: задайте \code{d = 1011} и по очереди \code{err} с одной единицей в каждой из семи позиций. На RTL найдите цепочки XOR (кодер и вычисление синдрома).

\stars{1} повторите: пройдите все семь позиций и запишите синдромы. \stars{2} найдите ошибку: подайте \code{err} с двумя единицами, убедитесь, что \code{q} теперь неверен, и объясните, почему код такого случая не лечит. \stars{3} проект: этап «Охранная панель» --- защита кода; критерий приёмки --- для всех 16 номеров зон и всех семи одиночных ошибок \code{q} совпадает с \code{d}.
\end{diybox}

\section{Три состояния, ШИМ и АЦП: что где разобрано}
\label{sec:b-analog-ref}

Узлы на границе кристалла и аналоговые узлы вынесены в главы, где разобрана их
теория, вместе с кодом, тестбенчами и заданиями. Здесь~--- указатель.

\begin{table}[H]
\centering
\caption{Узлы на границе кристалла: где разобраны и какой этап проекта закрывают}
\small
\begin{tabularx}{\textwidth}{l l X}
\toprule
\textbf{Узел} & \textbf{Где} & \textbf{Файлы и этап «Охранной панели»} \\
\midrule
Ножка с тремя состояниями, шина, открытый сток & гл.~\ref{ch:electrical} & \code{b5a\_electrical/tristate.v}, \code{tb\_bus.v}, \code{00\_theory/io\_pad.v}; сирена на общей линии: вне тревоги \code{oe = 0} \\
Подтяжка кнопки & гл.~\ref{ch:electrical} & \code{b5a\_electrical/btn\_in.v}, \code{PULL\_MODE} в \code{.cst} \\
ШИМ-ЦАП, зуммер & гл.~\ref{ch:analog} & \code{00\_theory/pwm\_dac.v}, \code{b9b\_analog/pwm8.v}, \code{buzzer.v}; в тревоге тон 2~кГц, 25\,\% \\
Сигма-дельта & гл.~\ref{ch:analog} & \code{b9b\_analog/sigma\_delta.v} \\
АЦП последовательного приближения & гл.~\ref{ch:analog} & \code{00\_theory/sar\_adc.v}; шум: \code{result[7:4] > 9} даёт \code{loud = 1} \\
Очередь FIFO (журнал тревог) & гл.~\ref{ch:fsmdesign} & \code{00\_theory/fifo8.v}, \code{b9a\_fsm/fifo\_ptr.v}; зоны 1, 3, 2 читаются в том же порядке \\
\bottomrule
\end{tabularx}
\end{table}

\section{Охранная панель на Verilog}
\label{sec:b-panel}

Вот все узлы вместе. Панель росла по этапам: датчик двери, охрана и ночной
режим, голосование «2 из 3», порог шума, запоминание тревоги, автомат режимов
и ввод кода, писк. Теперь каждому этапу соответствует несколько строк
Verilog. Проект состоит из трёх частей, и все они лежат в одном файле
\code{code/05a\_bridge/panel\_top.v}:
\begin{itemize}
  \item \textbf{ядро} \code{guard\_panel}: вся логика панели; оно ничего не знает
        о плате и целиком проверяется в симуляции;
  \item \textbf{фильтр входа} \code{in\_filter}: синхронизатор из двух триггеров
        и подавление дребезга (глава~\ref{ch:clocks});
  \item \textbf{верхний модуль} \code{panel\_top}: соединяет ядро с кнопками,
        выключателями, светодиодами и зуммером конкретной платы.
\end{itemize}
Такое разделение~--- главный приём практики: когда логика отделена от платы,
её можно проверить тестбенчем, а на другую плату перейти, сменив только верхний
модуль и файл ограничений.

\subsection{Ядро: вся логика панели}

Первая половина модуля~--- комбинационная часть. Параметры \code{TONE} и
\code{GAP} выбирают разряды счётчика для писка, чтобы в симуляции их можно было
уменьшить.

\begin{vcode}
module guard_panel #(
    parameter TONE = 13,         // разряд счётчика для тона (27 МГц: около 1,6 кГц)
    parameter GAP  = 22          // разряд для паузы между гудками (около 3,2 Гц)
) (
    input        clk, rst,                     // такт и синхронный сброс
    input        door, window, motion,         // датчики, 1 = сработал
    input        night,                        // ночной режим
    input  [3:0] noise,                        // уровень шума
    input        arm_btn,                      // «поставить на охрану»
    input        code_bit, code_step,          // ввод кода: бит и импульс «принять»
    output       alarm_led, armed_led,
    output       siren,
    output reg [1:0] state
);
    localparam OFF = 2'd0, ARMED = 2'd1, ALARM = 2'd2;

    // Этапы 2, 3, 6: комбинационная часть
    wire vote     = (door & window) | (door & motion) | (window & motion);
    wire loud     = noise > 4'd9;
    wire intruder = vote | loud | (door & night);

\end{vcode}

Вторая половина модуля: ввод кода, автомат и сирена.

\begin{vcode}
    // Этап 9: ввод кода 1-0-1 и фронт кнопки
    reg [2:0] hist;
    reg       arm_d;
    always @(posedge clk) begin
        arm_d <= arm_btn;
        if (rst)            hist <= 3'b000;
        else if (code_step) hist <= {hist[1:0], code_bit};
    end
    wire arm_pulse = arm_btn & ~arm_d;
    wire unlock    = code_step & ({hist[1:0], code_bit} == 3'b101);

    // Этапы 7 и 9: автомат OFF -> ARMED -> ALARM (ALARM запоминается до верного кода)
    always @(posedge clk) begin
        if (rst) state <= OFF;
        else case (state)
            OFF:     if (arm_pulse) state <= ARMED;
            ARMED:   if (unlock)        state <= OFF;
                     else if (intruder) state <= ALARM;
            ALARM:   if (unlock) state <= OFF;
            default: state <= OFF;
        endcase
    end

    // Этап 8: счётчик для писка (тон и пауза)
    reg [23:0] cnt;
    always @(posedge clk) cnt <= rst ? 24'd0 : cnt + 24'd1;

    assign alarm_led = (state == ALARM);
    assign armed_led = (state == ARMED);
    assign siren     = alarm_led & cnt[GAP] & cnt[TONE];
endmodule
\end{vcode}

Разберём, что откуда взялось. Строки с \code{wire}~--- комбинационная часть;
на RTL это компаратор и несколько вентилей. Блок \code{always @(posedge clk)} с
\code{hist}~--- сдвиговый регистр из трёх триггеров, ловящий код 1-0-1: код
считается введённым в такте, когда пришёл третий бит (\code{code\_step}).
Оператор \code{case}~--- автомат из трёх состояний; состояние ALARM
удерживается, пока не придёт верный код: это и есть память тревоги. Счётчик
\code{cnt} и два его разряда дают писк с паузами: разряд \code{k} переключается с
частотой $f/2^{k+1}$, поэтому при 27~МГц \code{cnt[13]} даёт около 1,6~кГц
(тон), а \code{cnt[22]}~--- около 3,2~Гц (пауза между гудками).

\subsection{Тестбенч: приёмка без рук}

Тестбенч проходит по всем этапам и печатает \code{guard\_panel: OK}, если
панель ведёт себя как задумано. Параметры \code{TONE = 1} и \code{GAP = 3}
сокращают писк до нескольких тактов, иначе пришлось бы ждать миллионы.

\vfile{\codepath/05a_bridge/tb_guard_panel.v}{code/05a\_bridge/tb\_guard\_panel.v}

\textbf{Как читается.} Сценарий соответствует пунктам приёмки: (1)~в состоянии
OFF датчики и шум не действуют; (2)~кнопка ставит на охрану, один датчик из
трёх тревоги не даёт; (3)~два датчика переводят в ALARM, закрытие двери тревогу
не снимает, а сирена то гудит, то молчит; (4)~неверная серия \code{1-1-0}
тревогу не снимает, верная \code{1-0-1} снимает; (5)~шум 9 не страшен, шум 10
включает тревогу, дверь ночью тоже. Задачи \code{tick}, \code{press\_arm},
\code{code} превращают их в одну строку каждая, поэтому сценарий читается как
описание.

Для запуска: \code{iverilog -g2012 -o tb tb\_guard\_panel.v panel\_top.v}, затем
\code{vvp tb}. В Logra тот же сценарий строят на вкладке «Диаграммы»: вставьте
\code{panel\_top.v}, выберите верхним \code{guard\_panel}, задайте параметры
\code{TONE = 1}, \code{GAP = 3} и подайте входы по таблице ниже.

\begin{table}[H]
\centering
\caption{Сценарий на вкладке «Диаграммы» (число циклов 120)}
\small
\begin{tabularx}{\textwidth}{l X l}
\toprule
\textbf{Шаг} & \textbf{Что подать} & \textbf{Что должно получиться} \\
\midrule
1 & импульс \code{rst} в начале; \code{door = 1}, \code{window = 1} & \code{state = 0}: в OFF датчики не действуют \\
2 & \code{arm\_btn = 1} на один такт & \code{state = 1}, горит \code{armed\_led} \\
3 & \code{door = 1}, \code{motion = 1} & \code{state = 2}, \code{siren} гудит; закрыть дверь~--- тревога держится \\
4 & \code{code\_bit = 1, 0, 1} с \code{code\_step = 1} на каждый бит & \code{state = 0}, \code{siren = 0} \\
\bottomrule
\end{tabularx}
\end{table}

\subsection{Привязка к плате}

Плата Sipeed Tang Nano 20K (ПЛИС GW2AR-LV18)~--- единственная плата, с которой
работают проекты Logra, поэтому привязка делается к ней. На плате есть кварц
27~МГц, две кнопки \code{S1}, \code{S2} (нажата~--- 1) и шесть светодиодов
(горят при 0). Датчиков и зуммера на плате нет: их подключают к ножкам
разъёма J5. Бит кода вводят кнопками: \code{S2}~--- единица, \code{S1}~---
нуль, каждое нажатие даёт импульс «принять».

\begin{table}[H]
\centering
\caption{Привязка портов \code{panel\_top} к плате Tang Nano 20K}
\small
\begin{tabular}{l l l l}
\toprule
\textbf{Порт} & \textbf{Что подключено} & \textbf{Вывод} & \textbf{Подтяжка} \\
\midrule
\code{clk} & кварц 27 МГц & 4 & --- \\
\code{btn[0]}, \code{btn[1]} & кнопки S1 (бит 0), S2 (бит 1) & 88, 87 & к земле \\
\code{door\_pin} & выключатель двери & 75 & к питанию \\
\code{window\_pin} & выключатель окна & 80 & к питанию \\
\code{motion\_pin} & датчик движения & 42 & к питанию \\
\code{night\_pin} & ночной режим & 41 & к питанию \\
\code{arm\_pin} & «поставить на охрану» & 56 & к питанию \\
\code{buzzer} & пассивный зуммер & 54 & --- \\
\code{led[0]}\ldots\code{led[5]} & светодиоды (0 = горит) & 15--20 & к питанию \\
\bottomrule
\end{tabular}
\end{table}

Выключатели замыкают ножку на землю, поэтому на ножке «не сработал» равно
единице (подтяжка к питанию), а внутри модуля сигнал инвертируется
(глава~\ref{ch:electrical}). Фильтр входа защищает панель от дребезга и
рассинхронизации; перед первым измерением он выдаёт «не сработал», чтобы при
включении не было ложной тревоги.

\begin{vcode}
// Вход с кнопки или датчика: два триггера (синхронизация) и фильтр дребезга.
module in_filter #(
    parameter N    = 270_000,    // сколько тактов значение должно держаться (10 мс при 27 МГц)
    parameter INIT = 1'b0        // значение «до первого измерения»
) (
    input      clk,
    input      raw,              // сигнал снаружи, не связанный с тактом
    output reg clean = INIT
);
    reg s1 = INIT, s2 = INIT;
    reg [18:0] cnt = 19'd0;
    always @(posedge clk) begin
        s1 <= raw;
        s2 <= s1;
        if (s2 == clean)          cnt <= 19'd0;
        else if (cnt == N - 1) begin clean <= s2; cnt <= 19'd0; end
        else                      cnt <= cnt + 19'd1;
    end
endmodule

\end{vcode}

Верхний модуль собирает всё вместе: сброс при включении, семь фильтров, ввод
кода двумя кнопками и ядро. Порты называются так же, как в файле ограничений.

\begin{vcode}
module panel_top (
    input        clk,            // кварц 27 МГц
    input  [1:0] btn,            // S1 = бит 0, S2 = бит 1 кода (нажата = 1)
    input        door_pin, window_pin, motion_pin, night_pin, arm_pin,   // выключатели на землю, разъём J5
    output       buzzer,
    output [5:0] led             // светодиоды горят при 0
);
    // Сброс при включении: 16 тактов после подачи питания
    reg [4:0] por = 5'd0;
    wire rst = ~por[4];
    always @(posedge clk) if (rst) por <= por + 5'd1;

    // Фильтры. Выключатели «на землю»: 1 на ножке = не сработал, поэтому INIT = 1 и инверсия.
    wire door_f, window_f, motion_f, night_f, arm_f, b0_f, b1_f;
    in_filter #(.INIT(1'b1)) f_door   (.clk(clk), .raw(door_pin),   .clean(door_f));
    in_filter #(.INIT(1'b1)) f_window (.clk(clk), .raw(window_pin), .clean(window_f));
    in_filter #(.INIT(1'b1)) f_motion (.clk(clk), .raw(motion_pin), .clean(motion_f));
    in_filter #(.INIT(1'b1)) f_night  (.clk(clk), .raw(night_pin),  .clean(night_f));
    in_filter #(.INIT(1'b1)) f_arm    (.clk(clk), .raw(arm_pin),    .clean(arm_f));
    in_filter                f_b0     (.clk(clk), .raw(btn[0]),     .clean(b0_f));
    in_filter                f_b1     (.clk(clk), .raw(btn[1]),     .clean(b1_f));

    // Ввод кода двумя кнопками: нажатие любой даёт импульс «принять», S2 --- бит 1, S1 --- бит 0
    reg any_d = 1'b0;
    always @(posedge clk) any_d <= b0_f | b1_f;
    wire code_step = (b0_f | b1_f) & ~any_d;
    wire code_bit  = b1_f;

    wire alarm_led, armed_led, siren;
    wire [1:0] state;
    guard_panel panel (
        .clk(clk), .rst(rst),
        .door(~door_f), .window(~window_f), .motion(~motion_f), .night(~night_f),
        .noise(4'd0),                          // датчика шума на плате нет (см. главу про АЦП)
        .arm_btn(~arm_f),
        .code_bit(code_bit), .code_step(code_step),
        .alarm_led(alarm_led), .armed_led(armed_led), .siren(siren), .state(state)
    );

    assign buzzer = siren;
    assign led    = ~{3'b000, state == 2'd0, armed_led, alarm_led};   // led[0] = тревога, led[1] = охрана, led[2] = выкл
endmodule
\end{vcode}

Файл ограничений связывает порты \code{panel\_top} с ножками: на каждый порт
две строки, \code{IO\_LOC} (номер вывода) и \code{IO\_PORT} (уровни, подтяжка,
ток).

\vfile[text]{\codepath/05a_bridge/panel_top.cst}{code/05a\_bridge/panel\_top.cst}

\textbf{Другая плата.} Ядро и тестбенч от платы не зависят. Чтобы перенести
проект, нужно (1) по описанию платы выписать выводы кварца, кнопок,
светодиодов и разъёма и составить новый \code{.cst}; (2) проверить полярность
светодиодов (у Tang Nano 20K горят при 0; если у другой платы при 1, уберите
инверсию в \code{assign led}); (3) пересчитать параметры по частоте кварца
$f$: \code{N} $\approx f\cdot 10$~мс (например, для 50~МГц это 500\,000, при
этом разрядность счётчика в \code{in\_filter} 19 бит достаточна до $5{,}2\cdot
10^5$), а разряды \code{TONE} и \code{GAP}~--- по формуле $f/2^{k+1}$ для тона
около 1,6~кГц и паузы около 3~Гц. На вкладке «Лаборатория» те же порты можно
подключить к виртуальным кнопкам, светодиодам и зуммеру до прошивки.

\begin{warnbox}
Сигналы кнопок и датчиков приходят снаружи и никак не связаны с тактом
панели. Без синхронизатора на плате панель время от времени будет «ошибаться»
(метастабильность, глава~\ref{ch:clocks}), а без фильтра дребезга одно
нажатие станет несколькими. Поэтому в \code{panel\_top} каждый внешний
сигнал проходит через \code{in\_filter}. В симуляции Logra и в тестбенче
сигналы идеальны, поэтому там ядро проверяют без фильтров.
\end{warnbox}

\begin{diybox}[Охранная панель: итог на Verilog и на плате]
\textbf{Задание.} Соберите \code{guard\_panel} из этой главы в Logra, проверьте все этапы тестбенчем и на диаграмме и подготовьте проект к плате.

\textbf{В Logra.} Проект «С нуля», вкладка «Код»: вставьте \code{panel\_top.v}, нажмите «Синтезировать» (Ctrl+Enter). Для диаграмм выберите верхним \code{guard\_panel} и задайте \code{TONE = 1}, \code{GAP = 3}, иначе писк придётся ждать миллионы тактов. Для платы верхним оставьте \code{panel\_top}, на вкладке «Выводы» назначьте порты по таблице главы (она повторяет файл \code{panel\_top.cst}), на вкладке «Подключение» прошейте плату.

\textbf{Ожидаемый результат.} После \code{arm\_btn} панель в состоянии ARMED. Срабатывание двух датчиков из трёх, шум выше 9 или дверь ночью переводят её в ALARM, и сирена гудит. Закрытие двери тревогу не снимает. Серия \code{1, 0, 1} на \code{code\_bit} с импульсами \code{code\_step} возвращает OFF. Тестбенч печатает \code{guard\_panel: OK}.

\textbf{Как проверить.} Диаграммы: сигналы \code{state}, \code{door}, \code{window}, \code{noise}, \code{alarm\_led}, \code{siren}; число циклов 120. RTL: найдите компаратор, три триггера \code{hist}, автомат. Вкладка Gowin покажет число ячеек. Для платы: после «Выводов» убедитесь, что порты \code{panel\_top} сопоставлены ножкам без предупреждений.

\stars{1} повторите: пройдите сценарий таблицы на диаграмме. \stars{2} допишите: добавьте выход \code{beep\_count}, который считает число тревог с начала работы (счётчик растёт на переходе в ALARM). \stars{3} Итоговая приёмка (этап «Охранная панель», плата): (1) панель встаёт на охрану по кнопке и не реагирует на датчики в состоянии OFF; (2) любая из трёх причин переводит её в ALARM, закрытие двери не гасит тревогу; (3) снять тревогу можно только кодом 1-0-1, любая другая серия не работает; (4) на плате Tang Nano 20K кнопки подключены через синхронизатор и антидребезг, зуммер даёт писк с паузами.
\end{diybox}

\textbf{Образец ответа к \stars{2}.} Тревогу считают по фронту: сравнивают текущее состояние с предыдущим.
\begin{vcode}
module alarm_counter (
    input        clk, rst,
    input        alarm_led,        // сигнал из guard_panel
    output reg [3:0] beep_count
);
    reg prev;
    always @(posedge clk) begin
        prev <= alarm_led;
        if (rst) beep_count <= 4'd0;
        else if (alarm_led && !prev) beep_count <= beep_count + 4'd1;
    end
endmodule
\end{vcode}


\begin{ideabox}
Весь «путь от идеи до схемы» повторяется каждый раз:
формулировка поведения $\to$ описание на Verilog $\to$ схема из вентилей и
триггеров $\to$ таблицы LUT и триггеры кристалла. Умение переходить между
этими уровнями и есть главный навык разработчика цифровых устройств.
\end{ideabox}

\begin{taskbox}[Задания]
\begin{enumerate}
  \item Опишите на Verilog 4-разрядный сумматор с входом переноса и выходом
        переполнения и проверьте его на вкладке «Диаграммы». Сколько элементов
        ALU использует синтезатор на вкладке Gowin?
  \item Добавьте в \code{counter4} вход \code{up}: при \code{up = 0} счётчик должен
        считать вниз. Что изменилось на схеме RTL?
  \item Измените светофор так, чтобы перед зелёным светом на два такта
        загорался вместе с красным жёлтый. Сколько состояний теперь нужно?
  \item Опишите мультиплексор $4:1$ оператором \code{assign y = d[s]} и
        сравните его схему на вкладке RTL с рисунком в главе «Мультиплексоры».
\end{enumerate}
\end{taskbox}
